\section{Taller resuelto: Qué parte de resistencia disminuye al consumir}

\noindent\textbf{Nivel:} Intermedio.

\apartado{1. Enunciado}{A TAS 98 m/s y densidad 0.92, un avión de 5100 kg con S=29 m² y polar 0.024+0.047CL² consume el 10 por ciento de su masa total. Calcule resistencia inicial y final manteniendo altitud y TAS.}

\apartado{2. Datos}{La masa final es 4590 kg; vuelo horizontal en ambos extremos.}

\apartado{3. Incógnitas}{Las magnitudes y comprobaciones solicitadas en el caso.}

\apartado{4. Sistema de referencia}{Fuerzas en ejes viento; TAS mantenida.}

\apartado{5. Hipótesis físicas}{Polar fija; se omite variación de resistencia de trim y de centro de gravedad.}

\apartado{6. Modelo matemático}{Qué parte de resistencia disminuye al consumir}

\apartado{7. Ecuaciones simbólicas}{Se parte de los balances y definiciones siguientes.}

\begin{equation}D=A_V+B_VW^2\end{equation}

\begin{equation}D_{parasita}=q_\infty SC_{D0}\end{equation}

\begin{equation}D_{cuadratica,f}=0.9^2D_{cuadratica,i}\end{equation}

\apartado{8. Desarrollo algebraico}{Las operaciones se realizan en este orden, conservando las condiciones del modelo.}

\begin{enumerate}

\item La presión dinámica y el término parásito permanecen constantes.

\item La sustentación adimensional disminuye en la proporción de masa.

\item El término cuadrático disminuye con el cuadrado de esa proporción, no linealmente.

\item Se suman ambas partes para comparar la resistencia total.

\end{enumerate}

\apartado{9. Sustitución numérica}{Se conservan los resultados intermedios sin redondear y se muestran todas las operaciones escalares disponibles.}

\textbf{Desarrollo de las magnitudes intermedias.} Con TAS y densidad fijas, $qS$ permanece constante. El peso cambia y obliga a cambiar $C_L=W/(qS)$. Por tanto el término parásito permanece fijo y el término cuadrático disminuye como el cuadrado de la masa.\begin{enumerate}

\item Calcular la fuerza de referencia constante.

\[\fitmath{\mathrm{qS}=0.5\cdot 0.92\cdot {98}^{2}\cdot 29\approx 128117.4\,\mathrm{N}}\]

\item Equilibrio con masa inicial.

\[\fitmath{\mathrm{CLi}=\frac{5100\cdot 9.80665}{\mathrm{qS}}}\]

\[\fitmath{\mathrm{CLi}=\frac{5100\cdot 9.80665}{(128117.36)}\approx 0.3903758\,}\]

\item Equilibrio con masa final.

\[\fitmath{\mathrm{CLf}=\frac{4590\cdot 9.80665}{\mathrm{qS}}}\]

\[\fitmath{\mathrm{CLf}=\frac{4590\cdot 9.80665}{(128117.36)}\approx 0.3513382\,}\]

\item Resistencia inicial de la polar.

\[\fitmath{\mathrm{Di}=\mathrm{qS}\cdot \left(0.024+0.047\cdot {\mathrm{CLi}}^{2}\right)}\]

\[\fitmath{\mathrm{Di}=(128117.36)\cdot \left(0.024+0.047\cdot {(0.3903757851)}^{2}\right)\approx 3992.455\,\mathrm{N}}\]

\item Resistencia final.

\[\fitmath{\mathrm{Df}=\mathrm{qS}\cdot \left(0.024+0.047\cdot {\mathrm{CLf}}^{2}\right)}\]

\[\fitmath{\mathrm{Df}=(128117.36)\cdot \left(0.024+0.047\cdot {(0.3513382066)}^{2}\right)\approx 3818.104\,\mathrm{N}}\]

\item Comparar fuerzas completas, no solo el término cuadrático.

\[\fitmath{\mathrm{reduction}=1-\frac{\mathrm{Df}}{\mathrm{Di}}}\]

\[\fitmath{\mathrm{reduction}=1-\frac{(3818.103746)}{(3992.455043)}\approx 0.0436702\,}\]

\end{enumerate}





\apartado{10. Resultados con unidades}{Los resultados pertenecen al modelo y al intervalo de datos declarados.}

\begin{center}\begin{tabular}{lr}\toprule Magnitud&Resultado\\\midrule

Resistencia inicial & $3992.455\,\mathrm N$\\

Resistencia final & $3818.104\,\mathrm N$\\

Reducción relativa & $0.0436702\,$\\

\bottomrule\end{tabular}\end{center}

\apartado{11. Interpretación}{Una reducción del 10 por ciento de masa no implica una reducción del 10 por ciento de resistencia total a TAS fija.}

\apartado{12. Comprobación independiente y alcance}{Calcular directamente CL inicial y final y evaluar la polar, contrastando con la separación $A_{V}$+$B_{V}$ W².}

\apartado{13. Limitaciones}{El factor cuadrático total puede contener resistencia de perfil dependiente de sustentación; no se lo denomina inducida pura sin justificación.}

\par\smallskip\noindent{\footnotesize\textit{Datos y solución originales. Las sustituciones registradas se reproducen con SymPy y Studium; eso no establece validación experimental ni revisión experta.}}
